\section*{Chapter \ref{ch:rc:portfolio-credit} --- Portfolio Credit}

\begin{solution}{exo:rc:portfolio-credit:1}
$125\times4.08\% = 5.10$ defaults. No: expectations are linear, so the expected number of
defaults depends only on the marginals; correlation changes the dispersion around it.
\end{solution}

\begin{solution}{exo:rc:portfolio-credit:2}
The equity tranche is long correlation: higher correlation raises the chance of few
defaults, which is where it survives. Senior tranches are short correlation: only
correlated scenarios reach them. The mezzanine sits between: at low correlation more
correlation pushes losses into it, at high correlation losses jump past it or none occur,
so its value is not monotone.
\end{solution}

\begin{solution}{exo:rc:portfolio-credit:3}
Each default costs $0.6/125 = 0.48\%$ of the pool: the equity is wiped out by the seventh
default ($3/0.48 = 6.25$), and losses reach 22\% at the forty-sixth ($22/0.48 = 45.8$).
\end{solution}

\begin{solution}{exo:rc:portfolio-credit:4}
\omterm{def:rc:portfolio-credit:dc}{Default correlation} is the correlation of indicators of rare events. Two latent variables
with correlation 25\% often move together, but both must fall below a threshold 1.74
standard deviations down for a joint default; the joint probability rises well above
$p^2$ but stays small relative to $p$, so the indicators' correlation is small.
\end{solution}

\begin{solution}{exo:rc:portfolio-credit:5}
$50\times3.28 = \text{EUR}~164$ million of bought index protection.
\end{solution}

\begin{solution}{exo:rc:portfolio-credit:6}
Base tranches' values are monotone in correlation, so a base correlation always exists and
is unique, and it interpolates across detachments; \omterm{def:rc:portfolio-credit:compound}{compound correlations} can be two or
none. But pricing a mezzanine as the difference of two base tranches at different
correlations is not a model: interpolated base correlations can give negative expected
losses for thin tranches, and bespoke pools and maturities need a mapping rule.
\end{solution}

\begin{solution}{exo:rc:portfolio-credit:7}
\texttt{tail\_table()}: 0.13\% under the \omterm{def:rc:portfolio-credit:gauss}{Gaussian copula} and 1.13\% under the Student-t with
four degrees of freedom, a ratio of 8.6.
\end{solution}

\begin{solution}{exo:rc:portfolio-credit:8}
Expected loss is not the risk: the value of a \omterm{def:rc:portfolio-credit:ss}{super senior tranche} moves with the
probability of extreme, correlated losses, which rises in a crisis exactly when the holder
cannot sell; its large notional turns small price moves into large losses; and a rating
based on a model's tail inherits the model's error there. It needs a hedge or capital for
its market value, as 2007--08 showed.
\end{solution}

\begin{solution}{pb:rc:portfolio-credit:1}
\textbf{1.} 17.3 and 7.32 times the index; EUR~23.56 million of 3--6\%.
\textbf{2.} EUR~3.721 million upfront; running, 500 basis points on the equity (EUR~500\,000 a
year) less 195 on the mezzanine (EUR~459\,502), about EUR~40\,000 a year while nothing
defaults.
\textbf{3.} Flat to a small parallel widening; long default risk of any single name (the
equity absorbs the first losses); long correlation.
\textbf{4.} It gains when correlation rises (the equity's value rises more than the
mezzanine's).
\textbf{5.} Idiosyncratic moves, correlation moves, the convexity of large moves, and the
change of the deltas over time.
\textbf{6.} Equity upfront 42.79\%; mezzanine 236.5 basis points.
\textbf{7.} Equity $-\text{EUR}~557\,577$, mezzanine $+\text{EUR}~424\,041$, total
$-\text{EUR}~133\,536$.
\textbf{8.} The delta is for all names moving together; one name's widening moves the
equity much more than the mezzanine, relative to their parallel deltas.
\textbf{9.} 0.48\% of the pool, 16\% of the equity tranche: EUR~1.6 million.
\textbf{10.} No: the index moves by one name's share, $1/125$, while the equity takes a large
part of that name's risk; single-name protection on the name is the hedge.
\textbf{11.} Equity $-\text{EUR}~269\,618$, mezzanine $-\text{EUR}~594\,873$ (its spread falls to 138
basis points), total $-\text{EUR}~864\,491$.
\textbf{12.} The 3--6\% expected loss is $2\,\mathrm{EL}_{0,6}-\mathrm{EL}_{0,3}$: with $\rho_6$ unchanged, a lower
$\rho_3$ raises $\mathrm{EL}_{0,3}$ and lowers the mezzanine's.
\textbf{13.} Equity $-\text{EUR}~828\,391$, mezzanine $-\text{EUR}~171\,869$, total $-\text{EUR}~1\,000\,261$.
\textbf{14.} The mezzanine protection, bought to gain on widening; its spread fell with
correlation, so it lost too.
\textbf{15.} Sellers of equity protection buy it back (widening the equity) and sell back their
mezzanine protection (tightening it), moving both legs against every similar position.
\textbf{16.} In the model, the asset correlation of the \omterm{def:rc:portfolio-credit:gauss}{Gaussian copula}; in the market, the
price of the equity tranche relative to the rest of the structure, set by supply and
demand.
\textbf{17.} Limits on correlation sensitivity and on single-name jump-to-default by name, stress
tests of joint idiosyncratic and correlation moves, and the liquidity of the hedges.
\textbf{18.} The \omterm{def:rc:portfolio-credit:gauss}{Gaussian copula} with a correlation fitted to calm markets had thin tails, and
the mortgage pools' correlations were far higher than assumed; the models measured
market prices, not the risk of the factor.
\textbf{19.} Named result: \emph{the correlation trade in May 2005}: the equity leg loses
EUR~828\,391 and the mezzanine hedge EUR~171\,869, a total of EUR~1\,000\,261 on EUR~10
million of equity.
\textbf{20.} On correlation and on the dispersion of names: the delta removes only parallel
spread moves.
\end{solution}

\begin{solution}{iq:rc:portfolio-credit:1}
$X_i = \sqrt\rho Z+\sqrt{1-\rho}\varepsilon_i$, default by $t$ if $X_i < N^{-1}(p_i(t))$;
conditional on $Z = z$, $p_i(t\mid z) = N((N^{-1}(p_i(t))-\sqrt\rho z)/\sqrt{1-\rho})$, and
names are independent.
\iqlookfor{the latent variable, the threshold and conditional independence.}
\end{solution}

\begin{solution}{iq:rc:portfolio-credit:2}
Condition on the factor; names are then independent, and the distribution of losses
follows by the recursion adding one name at a time (on a loss grid for unequal notionals);
integrate over the factor by Gauss--Hermite quadrature. Cost: names times loss grid times
quadrature nodes, per date.
\iqlookfor{conditional independence, recursion and quadrature.}
\end{solution}

\begin{solution}{iq:rc:portfolio-credit:3}
Compound: one correlation per tranche; mezzanines can have two or none. Base: the
correlation of the base tranche $[0,d]$, always unique; tranches priced as differences of
base tranches at different correlations, which can give arbitrage (negative expected
losses) and needs mapping for bespoke pools.
\iqlookfor{non-uniqueness versus inconsistency.}
\end{solution}

\begin{solution}{iq:rc:portfolio-credit:4}
By market value, not expected loss: sensitivities to spreads, correlation and the factor;
stress tests of systematic scenarios; tail models beyond the Gaussian; counterparty risk
of any protection bought (monolines, AIG); liquidity and concentration limits.
\iqlookfor{market-value risk, stress tests and hedge counterparties.}
\end{solution}

\begin{solution}{iq:rc:portfolio-credit:5}
The limit of the probability that one variable is extreme given that the other is. The
\omterm{def:rc:portfolio-credit:gauss}{Gaussian copula} has none (for correlation below one); the Student-t and Clayton copulas
have it. Senior tranches lose only when many names default together, which is exactly
what tail dependence governs.
\iqlookfor{the definition and the link to senior tranches.}
\end{solution}

\begin{solution}{iq:rc:portfolio-credit:6}
It was flat to parallel moves but long single-name risk and long correlation. Idiosyncratic
blow-outs (GM and Ford) widened the equity tranche far more than the mezzanine; the
correlation implied by equity prices fell, pushing the equity down further and the
mezzanine spread down, so the mezzanine protection lost too; unwinds amplified both.
\iqlookfor{idiosyncratic risk, correlation risk and crowding.}
\end{solution}
